\section{The complete critical family}\label{sec:critical}
The four even actions used in the model are listed below. The indicated
kernel is trivial in the first two cases and is the centre in the last
two. The resulting quotient is elementary abelian of the indicated rank.
The extraspecial action is the faithful transitive action of plus type
on eight points, denoted $E_8$.
\begin{table}[ht]
\centering
\begin{tabular}{lrrr}
\toprule
Action & Degree & Quotient rank & $a(U)$\\
\midrule
$C_2$ & 2 & 1 & 2\\
regular $V_4$ & 4 & 2 & 24\\
natural $D_8$ & 4 & 2 & 8\\
$E_8\cong 2_+^{1+4}$ & 8 & 4 & 384\\
\bottomrule
\end{tabular}
\caption{The critical even actions.}\label{tab:critical}
\end{table}
The normalizer and quadratic-form calculations are literal finite
calculations; Appendix~\ref{sec:certificates} describes their witnesses.
The theorem about this family does not require the additional assertion
that these exhaust equality in a permutation-rank bound.

For a fixed word $D$ in these actions, let $K$ be the product of the
specified kernels and put $V=D/K=\F_2^R$. If $t=\dim K$, the squaring
map is a vector-valued quadratic map $q:V\longrightarrow K$.
Its polarization is the commutator map. Write
\[
 \phi=\prod_{i\geq1}(1-2^{-i})>0.
\]
The Gaussian product formula implies
\begin{equation}\label{eq:gaussian-bounds}
 2^{k(R-k)}\leq\gauss Rk2\leq\phi^{-1}2^{k(R-k)},
 \qquad G_R\asymp2^{\lfloor R^2/4\rfloor}.
\end{equation}
All constants in this section are absolute.

\begin{lemma}[The canonical subgroups]\label{lem:canonical}
Uniformly in the critical word, the number of full-projecting subspaces
$U\leq V$ is $G_R(1+O(R2^{-R/2}))$.
\end{lemma}
\begin{proof}
Failure to project onto a quotient coordinate of dimension at most four
places $U$ in one of its at most $15$ pullback hyperplanes. There are at
most $R$ coordinates. Each hyperplane contains $G_{R-1}$ subspaces, and
\eqref{eq:gaussian-bounds} gives $G_{R-1}/G_R=O(2^{-R/2})$.
The full preimage of a full-projecting $U$ contains $K$. Its projection
to each factor therefore contains the specified kernel and surjects
onto the quotient, so equals that factor. Different $U$ give different
full preimages.
\end{proof}

\begin{lemma}[All lifts of a subspace]\label{lem:critical-lifts}
Let $U\leq V$ have dimension $k$, and set
$Q(U)=\langle q(u):u\in U\rangle\leq K$. Put
$\mathcal A(U)=Q(U)^\perp\leq K^*$. The number of subgroups of $D$
whose image is exactly $U$ is
\begin{equation}\label{eq:critical-lifts}
 \sum_{B\leq\mathcal A(U)}2^{k\dim B}
 =\sum_{j=0}^{\dim\mathcal A(U)}
       \gauss{\dim\mathcal A(U)}j2\,2^{kj}.
\end{equation}
If $U$ is full-projecting, all these subgroups project onto the factors.
\end{lemma}
\begin{proof}
For such a subgroup $H$, put $W=H\cap K$. Necessarily $q(U)\subseteq W$.
Conversely every element of $\pi^{-1}(U)/W$ then has square one.
Such a group is abelian, since $(xy)^{-1}=y^{-1}x^{-1}$ and every
element is its own inverse; hence the quotient is elementary abelian.
The subgroups with intersection $W$ are the complements to $K/W$ in
this elementary abelian extension. They form a torsor under
$\Hom(U,K/W)$ and number $2^{k\codim_K W}$. Annihilator duality
$B=W^\perp$ gives \eqref{eq:critical-lifts}.
Fullness follows again from Burnside's basis theorem on each factor.
\end{proof}

\begin{lemma}[Quadratic realization]\label{lem:quadratic}
Fix a quadratic form $Q$ on a binary vector space. If $q_0$ is a
nonsingular plus-type form of dimension $h\in\{2,4\}$, the surjections
$f$ satisfying $q_0\circ f=Q$ form either the empty set or a torsor for
$O^+(h,2)$. In particular there are at most $72$ such maps.
\end{lemma}
\begin{proof}
If a map exists, the radical of the polar form of $Q$ is exactly
$\ker f$, since the target polar form is nonsingular and $f$ is onto.
Thus every such map has the same kernel and induces an isometry from
the same quotient to the target. Any two differ by a unique target
isometry. The elementary orders are $|O^+(2,2)|=2$ and
$|O^+(4,2)|=72$.
\end{proof}

\begin{proposition}[Uniform exceptional-lift estimate]\label{prop:extra-lifts}
Set $d=R/2-t\geq0$. For the fixed critical word, the total number of
noncanonical full subgroups is $O(G_R2^{-d})$.
\end{proposition}
\begin{proof}
Fix an $\ell$-space $B\leq K^*$, where $\ell\geq1$, and put its
equations in row-echelon form in the coordinates of the nonabelian
critical factors. Count ordered maps $\F_2^k\to V$ which are injective
and full on each coordinate and which satisfy $B\leq\mathcal A(U)$.
Once the nonpivot coordinate maps have been chosen, each pivot equation
prescribes one quadratic form. Lemma~\ref{lem:quadratic} allows at most
$72$ choices for its surjective coordinate map. Each pivot has target
dimension at least two. Thus the ordered-map count is at most
$72^\ell2^{(R-2\ell)k}$.

Division by $|\GL(k,2)|$, compared with the unrestricted ordered
injection count at least $\phi2^{Rk}$, gives
\[
 \#\{U:\dim U=k,\ B\leq\mathcal A(U),\ U\text{ full}\}
 \leq\phi^{-1}\gauss Rk2\,72^\ell2^{-2\ell k}.
\]
There are at most $\phi^{-1}2^{\ell(t-\ell)}$ choices for $B$.
Multiplying by its lift weight $2^{k\ell}$ in
\eqref{eq:critical-lifts}, the exceptional mass at dimension $k$ is at
most
\begin{equation}\label{eq:exception-k}
 \phi^{-2}\gauss Rk2\sum_{\ell\geq1}
 72^\ell2^{-\ell(k-t+\ell)}.
\end{equation}
Put $x=k-R/2$. Dividing by $G_R$ introduces a factor at most
$C2^{-x^2}$ by \eqref{eq:gaussian-bounds}. The exponent depending on
$x$ and $\ell$ becomes
\[
 -x^2-\ell x-\ell d-\ell^2
 =-(x+\ell/2)^2-\ell d-3\ell^2/4.
\]
The Gaussian sum over either lattice is uniformly bounded. Also,
\[
 \sum_{\ell\geq1}72^\ell2^{-3\ell^2/4}<\infty.
\]
Since $d\geq0$, the result is $O(2^{-d})$.
\end{proof}

\begin{proof}[Proof of Theorem~\ref{thm:critical-intro}]
Summing \eqref{eq:weights} over critical profiles of rank $R$ gives
$n!c_R$ in even degree, because the two rank-two actions have combined
weight $1/24+1/8=1/6$. The canonical count therefore supplies the main
term, with the uniform error of Lemma~\ref{lem:canonical}.

For multiplicities $(a,v,b,e)$ of the four actions,
\[
 R=a+2v+2b+4e,\quad t=b+e,\quad d=a/2+v+e.
\]
Inserting the weight $2^{-d}$ of Proposition~\ref{prop:extra-lifts}
changes the exponential generating polynomial to
\begin{equation}\label{eq:tilde-P}
 \widetilde P(y)=\frac{y}{2\sqrt2}+\frac{7y^2}{48}
                         +\frac{y^4}{768}.
\end{equation}
For fixed $a_1,a_2,a_4>0$ one has
\begin{equation}\label{eq:quartic-coefficient}
 \log_e[y^R]\exp(a_1y+a_2y^2+a_4y^4)
 =-\frac R4\log_e R+\frac R4(1+\log_e(4a_4))+O(\sqrt R).
\end{equation}
For completeness, the upper bound follows by substituting
$y=(R/(4a_4))^{1/4}$ into the positive series. For the lower bound take
$\lfloor R/4\rfloor$ quartic terms and at most three linear terms;
Stirling's inequalities give the same two terms with error $O(\log R)$.
Halving the quartic coefficient in \eqref{eq:tilde-P} therefore gives
exceptional relative mass $2^{-R/4+O(\sqrt R)}$.

In odd degree the fixed-point option contributes $c_R$. The natural
$S_3$ option contributes $c_{R-1}/6$. Any subgroup full on its $S_3$
coordinate and on the binary coordinates contains $A_3$ on the first
coordinate alone: raise an element projecting to a $3$-cycle to a
sufficiently large power of two. Quotienting this forced subgroup turns
that coordinate into one elementary binary coordinate. Its full lift
analysis is consequently the same one, with exceptional marker
$y/(6\sqrt2)$. Formula~\eqref{eq:quartic-coefficient}, also after
a fixed coefficient shift, gives the same exponential conclusion.
The two options have different actual orbit profiles and are disjoint.
\end{proof}
